\begin{description}
\item[(A)] Use the relation between apparent and absolute magnitude,
\[
m-M=5\log d-5.
\]
The distance from the parallax is $d=1/0.350=2.86\,\mathrm{pc}$, so Sirius A has
\[
M_V=-1.47-5\log2.86+5=1.25\ \mathrm{mag}.
\]
Taking the bolometric correction into account gives $M_B=1.25-0.4=0.85$ mag.
For Sirius B the same process gives $M_B=0.85+(8.5+1.47)=10.82$.
Use
\[
M_V=-2.5\log\frac{L}{L_0},\qquad L=L_0\times10^{-0.4M_V}.
\]
% sic: the official solution labels the bolometric magnitude M_V in this luminosity relation.
Substitution gives $L_A=1.37\times10^{28}\,\mathrm{W}$ and $L_B=1.41\times10^{24}\,\mathrm{W}$.
The Stefan--Boltzmann law $L=4\pi R^2\sigma T^4$ then gives
$R_A=1.39\times10^9\,\mathrm{m}$ and $R_B=2.25\times10^6\,\mathrm{m}$. \hfill\mbox{(12 points)}
\item[(B)] Convert the angular separation to radians,
\[
\theta_1=\frac{7}{206265}=3.39\times10^{-5}\ \mathrm{rad}.
\]
The separation is $s_1=3.39\times10^{-5}\times2.86=9.71\times10^{-5}\,\mathrm{pc}$,
or $3.00\times10^{12}\,\mathrm{m}$. \hfill\mbox{(8 points)}
\item[(C)] Use $M_{\mathrm{tot}}=a^3/P^2$, with $a$ in astronomical units, $P$ in years,
and mass in solar masses. The separation is 20 au, giving
\[
M_{\mathrm{tot}}=\frac{20.0^3}{50^2}=3.20\,M_\odot.
\]
\hfill\mbox{(8 points)}
\item[(D)] Sirius B has mass $\tfrac13\times3.20=1.07\,M_\odot$ and radius
$2.25\times10^6\,\mathrm{m}$ from part (A). Its mean density is
$\bar\rho=4.4\times10^{10}\,\mathrm{kg\,m^{-3}}$. \hfill\mbox{(6 points)}
\item[(E)] Its radius is comparable to Earth's while its mass is close to the Sun's.
It is therefore a white dwarf, closer to star 1. \hfill\mbox{(6 points)}
\end{description}
